% Revised Discussion paragraphs (2026-09-02). Every number below was checked
% against results/ ; see comments for the source file of each figure.

\paragraph{The format trap is a property of the training set, not of the data source.}
Devbunova's probes could have followed format because her evaluation and
deployment prompts came from different corpora. Our design removes that
explanation: the same payload appears in every cell, and the confounded probe
still classifies by format at every layer of every checkpoint. The depth
profiles show why. Format is decodable at layer 0 (AUC $\geq$ .999), before
any computation has occurred, whereas purpose only forms by 10--20\% depth. A
probe whose labels co-vary with format is offered a strong, early, trivially
separable feature and a weaker, later one, and it takes the former. Any
evaluation-awareness probe trained on natural data, where evaluation prompts
are overwhelmingly benchmark-shaped, faces the same choice.

% Sources: results/layer_sweep_full320-v3_*.json (format AUC >= .9998 at L0;
% decorrelated AUC crosses .98 at L6/32, L8/80, L8/80);
% results/full320_v3_matched_format_projection_diagnostic.json
% (rescue_fraction_median_excl_near_zero = .720/.549/.743;
%  mean_abs_delta_joint_auc = .0013/.0006/.0007; per-layer max .0029).
\paragraph{Stated purpose is largely, but not entirely, format-general.}
With format balanced, purpose is decodable at .989--.998 AUC, and a probe that
has seen only one format loses at most .05 AUC on the other, so the purpose
direction is mostly shared across formats. The residual gap is itself
informative. In a preliminary analysis, we estimate a rank-64 format subspace
from the matched benchmark--casual contrast, project it out, and retrain the
probes. Across layers where purpose has formed (all-format AUC $\geq$ .95),
this removes a median 55--74\% of the cross-format gap in all three checkpoints
while changing the all-format AUC by a mean absolute .0006--.0013 (never more
than .003 at any layer). Format and purpose therefore occupy largely separable
subspaces, and most of the transfer penalty lives in the format subspace rather
than in a format-specific encoding of purpose.

% Sources: results/nocue_sweep_full320-v3_*.json at L8/L18/L18/L37
% (format AUC .469/.731/.697/.672; 405B pct_eval 10.91 / 6.25; 405B
% decorrelated AUC .9980). 405B is an unpinned NDIF revision (README).
\paragraph{Balanced training does not yield a format-blind probe.}
Decorrelation removes the confound from the training set without removing
format from the probe: at 70B, benchmark formatting alone moves a purpose-free
prompt toward ``evaluation'' (format AUC .70--.73), and the residual varies
with depth, so a single layer can under- or overstate it. We have also run the
sweep on Llama-3.1-405B-Instruct (an unpinned revision, reported as
development-tier evidence); at its selected layer (37 of 126) the decorrelated
probe reaches .998 AUC yet calls 10.9\% of no-cue benchmark and 6.2\% of casual
prompts ``evaluation'' (format AUC .67). We therefore suggest reporting no-cue
format AUC as a standard diagnostic alongside probe accuracy: a probe can be
near-perfect on labelled data and still respond to format on prompts that say
nothing about purpose.

% Sources: results/steering/analysis_{8b,70b,405b}_logit.json, strength_z vs
% five random directions (>=2 strong, >=1 suggestive, else null):
%   fmt      .31 null | 2.21 strong | 1.57 suggestive
%   pur_dec  .19 null |  .62 null   | 1.31 suggestive
%   pur_conf .28 null | 2.02 strong | 1.74 suggestive
% Models steered: Llama-3.1-8B L8, Llama-3.1-70B L18, Llama-3.1-405B L37.
% Llama-3.3-70B was never steered. Readout is eval log-odds, not a rate.
\paragraph{Decodability is not used.}
The results above are observational. In preliminary steering experiments on
the Llama-3.1 checkpoints (8B, 70B, and 405B), we add each probe direction to
the residual stream at the selected layer, at a matched norm, on no-cue
prompts, and compare the resulting shift in a verbalized evaluation/deployment
log-odds readout against random directions of the same norm. At Llama-3.1-70B
the format direction and the format-aligned (confounded) purpose direction both
move the readout well outside the random band, whereas the format-orthogonal
purpose direction stays inside it. At 8B no direction, including format, leaves
the band. At 405B every direction, including the format-orthogonal one, falls
between one and two random-direction standard deviations, so the null for the
orthogonal component is established only at 70B. Llama-3.3-70B was not steered.
So far the only clearly active component is the format-aligned one;
decodability of purpose should not be read as evidence that the model uses it.

% Sources: results/devbunova_decomp/decomp_{8b,70b,405b}.json (template regime,
% L8/L18/L37): confounded pct_eval 92.2/94.6/95.2; decorrelated 6.5/4.1/13.2;
% decorrelated within_bench AUC .589/.638/.583; overall AUC .580/.269/.354.
% Llama-3.3-70B has no Devbunova extraction (docs/devbunova_decomposition_plan.md).
% results/implicit_v2_{8b,llama31_70b_l18,llama33_70b_l18}_analysis.json:
% joint_format_auc .598 [.527,.670] / .638 / .660.
\paragraph{What is decoded is the stated purpose.}
Our labels are explicit statements, so a positive result shows that the model
represents what it was told, not that it infers evaluation status. Two
preliminary checks make the scope concrete. Applying our decorrelated probe to
Devbunova's items on the Llama-3.1 checkpoints largely dissolves her trap
(4.1\%, 6.5\%, and 13.2\% of benchmark-format deployment prompts called
``evaluation'' at 70B, 8B, and 405B, versus 92--95\% for the confounded probe),
yet within the benchmark format its AUC on her evaluation/deployment labels is
only .58--.64, and pooled across her formats it is at or below chance, since
those labels are never stated in her text. On scenarios where evaluative
structure is described indirectly rather than announced, the frozen probes
transfer at .598--.660 AUC, with the 8B interval reaching down to .53. The
stated-purpose signal is clean and format-general, but it decays quickly as
purpose becomes implicit.

\paragraph{Limitations.}
Results cover one model family, 16 topics, and synthetic payloads; layers are
chosen by held-out AUC, which flatters selected-layer values but not the depth
curves; and no behavioural consequence of the probed signal is established. The
preliminary analyses are development-tier and have not been run under sealed
protocols; the steering and Devbunova analyses cover the Llama-3.1 checkpoints
only, and the 405B checkpoint is an unpinned revision.
